\bookchapter{附錄A：新手機開箱“保姆級”流程}

\booksection{一、開箱驗機“三部曲”：}

\begin{itemize}
  \item \textbf{觀外觀}：檢查機身、屏幕有無劃痕、磕碰。
  \item \textbf{對信息}：覈對盒上之IMEI/序列號與發票上者為否一致。
  \item \textbf{試功能}：撥一電話試聽筒話筒，拍一張照檢查攝像頭。
\end{itemize}

\booksection{二、激活與初始設置“重點提醒”：}

\begin{itemize}
  \item \textbf{網絡}：先連上穩定之Wi-Fi。
  \item \textbf{賬號}：登錄君之蘋果ID、華爲賬號、小米賬號等。
  \item \textbf{基礎設置}：將"查尋我之手機"打開，設一靠譜之鎖屏密碼。
  \item \textbf{eSIM設置}：若君之手機支持eSIM（如iPhone、部分安卓旗艦），可於設置中添加eSIM套餐，無需插入實體SIM卡。出國時eSIM尤便，可徑購當地運營商之eSIM套餐。
  \item \textbf{數據遷移}：從舊手機遷移數據時，確保兩臺手機電量充足。iPhone用"快速開始"，安卓用官方換機App。
\end{itemize}

\booksection{三、裝App“安全第一”：}

\begin{itemize}
  \item \textbf{首選官方商店}：蘋果用App Store，安卓用手機自帶之應用商店。
  \item \textbf{權限勿亂給}：安裝時彈出之權限請求，想清楚"此App真的需要嗎？"
\end{itemize}
